% Junior Level 2, Week 9 -- answer key
\answersweek{9}{Project Studio}{Answers for workbook chapter 9 and play-at-home card 9.}

\section{Lesson questions}
\begin{answer}[Exit question]
Any must-have named with a real test is right, for example: \enquote{The forever: I clicked the flag and the bananas kept falling}, \enquote{The if touching: I caught a banana and heard the pop}, \enquote{The ending: I missed three on purpose and it said Game over}. A child who says \enquote{I built it all and then tested} has not yet got the one-at-a-time habit: ask them to show one test now.
\end{answer}

\begin{answer}[Reflection]
Any honest answer is right. Five colored stars means every must-have works (milestone M2). Fewer: plan extra studio time (buffer B2, or 15 minutes at the start of week 10). What the Player found is often the most interesting answer: write it down for the week 10 checklist row \enquote{playtests and fixes a bug a player found}.
\end{answer}

\begin{answer}[Ask before you share (safety talk)]
Leo asks Priya first: \enquote{May I show your maze to my cousin?} If she says no, he does not show it. Games stay with our family; nothing is put on the internet or sent to people we do not know.
\end{answer}

\section{Workbook exercises}
\begin{answer}[Exercise 9.1: My Game Plan Card (milestone M1)]
Every card is different. Look for: catch game or maze game chosen; a player; something to catch or to avoid; a way to win that can be reached (for example 10 catches, or touching the goal); a way to lose (for example 3 misses); one random surprise. In the paper prototype the partner could win and could lose. Too many ideas: keep the must-haves, put the extras on the back.
\end{answer}

\begin{answer}[Exercise 9.2: Build My Game (milestone M2)]
\textbf{What you should see in Mei's sample:} the bananas start at the hidden Star, fall, and wait 1 to 3 seconds at the top before each new fall. A catch adds 1 to the score with a pop; a banana that reaches the bottom adds 1 to misses. The Bowl's referee says the goal, \enquote{Catch 10 bananas!}, at the start, \enquote{You win!} at 10 and \enquote{Game over!} at 3 misses, then everything stops.

\textbf{What to look for in every game:} \tickbox{} a forever loop \quad \tickbox{} an if-touching check \quad \tickbox{} a score or misses variable that starts at 0 \quad \tickbox{} random (a pick random or go to random position) \quad \tickbox{} a win or game-over ending with stop all \quad \tickbox{} each must-have tested before the next one. Maze games: the player is sent back to the Star when it touches the walls (a miss) and wins on touching the goal. Full reference for both kinds: \texttt{exercises/week09/ex-build-my-game/solution.md} and the Project Teacher Guide.
\end{answer}

\begin{answer}[Exercise 9.3: Playtest]
Any honest list is right; one bug and one fix is enough. Common finds and fixes: the score does not go back to 0 (set it to 0 at the flag); misses count at the start (drag the Star marker away from the edge); an if outside the forever (drag it in); the Player did not know how to win (say the goal at the start); far too fast (a smaller number in \enquote{move}); the thing goes through the player (the if touching is missing or names the wrong sprite). A \enquote{not yet} is a good find, not a failure.
\end{answer}

\section{Play-at-home card}
\begin{answer}[Practice the showcase]
Look for the game's name, who the player is, how to win and how to lose, and one design choice with a reason (\enquote{I chose 3 misses, because 1 was too hard for my little sister}). A choice that came from a Player's wish is a great answer.
\end{answer}

\begin{answer}[Arcade poster]
Any poster with the game's name in big letters, a drawing of the player and the thing to catch or avoid, and how to win is right. Look for the games-for-everyone ideas from week 7: big words, and shapes as well as colors.
\end{answer}
